\begin{tomtat}
	\begin{minipage}[t]{0.3\textwidth}
		\vspace{0pt}
		\begin{tomtat}
			\begin{center}
				\textbf{ĐỀ 34}
			\end{center}
			
		\end{tomtat}
	\end{minipage}
	\hfill
	\begin{minipage}[t]{0.5\textwidth}
		\vspace{0pt}
		\begin{tomtat}
			\begin{center}
				\includegraphics[scale=0.08]{images/34}
			\end{center}
		\end{tomtat}
	\end{minipage}
\end{tomtat}
\textbf{I. PHẦN TRẮC NGHIỆM}
\begin{ex}
	Với giá trị nào của x thì biểu thức $A=\sqrt{2-3x}$ xác định?
	\choice{$x \ge \frac{2}{3}$}{$x \le \frac{2}{3}$}{$x \le \frac{3}{2}$}{$x \ge \frac{3}{2}$}
	\loigiai{
		Biểu thức $A = \sqrt{2-3x}$ xác định khi $2-3x \ge 0$.
		\begin{align*} 2-3x &\ge 0 \\ -3x &\ge -2 \\ x &\le \frac{-2}{-3} \\ x &\le \frac{2}{3} \end{align*}
		\\ \textbf{Đáp Án B}
	}
\end{ex}

\begin{ex}
	Nghiệm của phương trình $(2x-3)(x+2)=0$ là:
	\choice{$x=\frac{3}{2}; x=2$}{$x=-\frac{3}{2}; x=2$}{$x=\frac{3}{2}; x=-2$}{$x=-\frac{3}{2}; x=-2$}
	\loigiai{
		Ta có: $(2x-3)(x+2)=0$
		$\Leftrightarrow 2x-3=0$ hoặc $x+2=0$
		Trường hợp 1: $2x-3=0 \Leftrightarrow 2x=3 \Leftrightarrow x=\frac{3}{2}$.
		Trường hợp 2: $x+2=0 \Leftrightarrow x=-2$.
		Vậy nghiệm của phương trình là $x=\frac{3}{2}$ hoặc $x=-2$.
		\\ \textbf{Đáp Án C}
	}
\end{ex}

\begin{ex}
	Cho tam giác ABC vuông tại A có $AB=6$ cm, $\tan B = \frac{2}{3}$. Độ dài cạnh AC là:
	\choice{$4 \text{ cm}$}{$9 \text{ cm}$}{$2\sqrt{5} \text{ cm}$}{$3\sqrt{5} \text{ cm}$}
	\loigiai{
		Trong tam giác ABC vuông tại A:
		$\tan B = \frac{\text{cạnh đối}}{\text{cạnh kề}} = \frac{AC}{AB}$.
		Theo đề bài, $\tan B = \frac{2}{3}$ và $AB=6$ cm.
		$\frac{AC}{AB} = \frac{2}{3}$
		$\frac{AC}{6} = \frac{2}{3}$
		$AC = 6 \cdot \frac{2}{3} = 4$ (cm).
		\\ \textbf{Đáp Án A}
	}
\end{ex}

\begin{ex}
	Biết đồ thị hàm số $y=ax^2 (a \neq 0)$ đi qua điểm $M(-2;-2)$. Giá trị của a là:
	\choice{$1$}{$-1$}{$\frac{1}{2}$}{$-\frac{1}{2}$}
	\loigiai{
		Đồ thị hàm số $y=ax^2$ đi qua điểm $M(-2;-2)$, nên tọa độ của M thỏa mãn phương trình hàm số:
		$-2 = a(-2)^2$
		$-2 = a \cdot 4$
		$a = \frac{-2}{4}$
		$a = -\frac{1}{2}$.
		\\ \textbf{Đáp Án D}
	}
\end{ex}

\begin{ex}
	Cho bảng tần số ghép nhóm về thời gian (phút) đi từ nhà đến trường của một nhóm học sinh như sau:
	\begin{center}
		\begin{tabular}{|c|c|c|c|c|}
			\hline
			Nhóm & [7;13) & [13;19) & [19;25) & [25;31) \\
			\hline
			Tần số & 5 & 10 & 20 & 15 \\
			\hline
		\end{tabular}
	\end{center}
	Tần số tương đối ghép nhóm của nhóm [13;19) là:
	\choice{$10\%$}{$20\%$}{$30\%$}{$40\%$}
	\loigiai{
		Tổng số học sinh (tổng tần số) là $N = 5+10+20+15 = 50$.
		Tần số của nhóm [13;19) là $n_{[13;19)} = 10$.
		Tần số tương đối của nhóm [13;19) được tính bằng công thức:
		$f_{[13;19)} = \frac{n_{[13;19)}}{N} \times 100\%$.
		$f_{[13;19)} = \frac{10}{50} \times 100\% = \frac{1}{5} \times 100\% = 0,2 \times 100\% = 20\%$.
		\\ \textbf{Đáp Án B}
	}
\end{ex}

\begin{ex}
	Cặp số $(4;2)$ là nghiệm của hệ phương trình nào sau đây?
	\choice{$\begin{cases} x+y=6 \\ x-y=2 \end{cases}$}{$\begin{cases} x+y=6 \\ 2x-5y=8 \end{cases}$}{$\begin{cases} x-y=2 \\ x-5y=4 \end{cases}$}{$\begin{cases} x-y=2 \\ 3x+2y=5 \end{cases}$}
	\loigiai{
		Thay $x=4, y=2$ vào từng hệ phương trình:
		A. $\begin{cases} 4+2=6 \text{ (Đúng)} \\ 4-2=2 \text{ (Đúng)} \end{cases}$. Vậy $(4;2)$ là nghiệm của hệ này.
		B. $\begin{cases} 4+2=6 \text{ (Đúng)} \\ 2(4)-5(2)=8-10=-2 \neq 8 \text{ (Sai)} \end{cases}$.
		C. $\begin{cases} 4-2=2 \text{ (Đúng)} \\ 4-5(2)=4-10=-6 \neq 4 \text{ (Sai)} \end{cases}$.
		D. $\begin{cases} 4-2=2 \text{ (Đúng)} \\ 3(4)+2(2)=12+4=16 \neq 5 \text{ (Sai)} \end{cases}$.
		\\ \textbf{Đáp Án A}
	}
\end{ex}
\begin{ex}
	Đội Sao đỏ của trường THCS có 10 học sinh gồm 6 học sinh nam và 4 học sinh nữ. Chọn ngẫu nhiên 1 học sinh từ đội Sao đỏ để chấm thi đua. Xác suất để chọn được học sinh nữ bằng:
	\choice{$\frac{1}{6}$}{$\frac{1}{10}$}{$\frac{2}{5}$}{$\frac{3}{5}$}
	\loigiai{
		Tổng số học sinh trong đội Sao đỏ là 10.
		Số học sinh nữ là 4.
		Xác suất để chọn được một học sinh nữ là tỉ số giữa số học sinh nữ và tổng số học sinh:
		$P(\text{chọn nữ}) = \frac{\text{Số học sinh nữ}}{\text{Tổng số học sinh}} = \frac{4}{10} = \frac{2}{5}$.
		\\ \textbf{Đáp Án C}
	}
\end{ex}

\begin{ex}
	Cho tam giác ABC vuông tại A. AB=3cm; AC=4cm. $\sin B = ?$
	\choice{$\frac{4}{3}$}{$\frac{3}{4}$}{$\frac{4}{5}$}{$\frac{3}{5}$}
	\loigiai{
		Trong tam giác ABC vuông tại A, để tính $\sin B$, ta cần biết cạnh đối AC và cạnh huyền BC.
		Cạnh đối AC = 4cm.
		Áp dụng định lý Pitago để tính cạnh huyền BC:
		$BC^2 = AB^2 + AC^2 = 3^2 + 4^2 = 9 + 16 = 25$.
		$BC = \sqrt{25} = 5$ cm.
		Vậy, $\sin B = \frac{\text{cạnh đối}}{\text{cạnh huyền}} = \frac{AC}{BC} = \frac{4}{5}$.
		\\ \textbf{Đáp Án C}
	}
\end{ex}

\begin{ex}
	Hộp sữa có dạng hình trụ với đường kính đáy là 12cm, chiều cao là 18cm. Thể tích của hộp sữa đó bằng?
	\choice{$648\pi \text{ cm}^3$}{$432\pi \text{ cm}^3$}{$216\pi \text{ cm}^3$}{$2592\pi \text{ cm}^3$}
	\loigiai{
		Đường kính đáy của hình trụ là $d = 12$ cm.
		Bán kính đáy của hình trụ là $r = \frac{d}{2} = \frac{12}{2} = 6$ cm.
		Chiều cao của hình trụ là $h = 18$ cm.
		Thể tích của hình trụ được tính bằng công thức $V = \pi r^2 h$.
		$V = \pi \cdot (6)^2 \cdot 18$
		$V = \pi \cdot 36 \cdot 18$
		$V = 648\pi \text{ (cm}^3\text{)}$.
		\\ \textbf{Đáp Án A}
	}
\end{ex}
\newpage
\textbf{II. PHẦN TỰ LUẬN}
\begin{ex}
	Rút gọn biểu thức A biết $A = \frac{\sqrt{x}+1}{x+4\sqrt{x}+4} : \left( \frac{x}{x+2\sqrt{x}} + \frac{x}{\sqrt{x}+2} \right)$, với $x > 0$.
	\loigiai{
		Điều kiện: $x > 0$.
		Ta có $x+4\sqrt{x}+4 = (\sqrt{x}+2)^2$.
		Và $x+2\sqrt{x} = \sqrt{x}(\sqrt{x}+2)$.
		\begin{align*} A &= \frac{\sqrt{x}+1}{(\sqrt{x}+2)^2} : \left( \frac{x}{\sqrt{x}(\sqrt{x}+2)} + \frac{x\sqrt{x}}{\sqrt{x}(\sqrt{x}+2)} \right) \\ &= \frac{\sqrt{x}+1}{(\sqrt{x}+2)^2} : \left( \frac{\sqrt{x}}{\sqrt{x}+2} + \frac{x\sqrt{x}}{\sqrt{x}(\sqrt{x}+2)} \right) \quad (\text{Rút gọn } \frac{x}{\sqrt{x}(\sqrt{x}+2)} = \frac{\sqrt{x}}{\sqrt{x}+2}) \\ &= \frac{\sqrt{x}+1}{(\sqrt{x}+2)^2} : \left( \frac{\sqrt{x}}{\sqrt{x}(\sqrt{x}+2)} + \frac{x}{(\sqrt{x}+2)\sqrt{x}} \right) \quad (\text{Sai ở bước trên, làm lại phần trong ngoặc thứ hai}) \end{align*}
		Làm lại phần trong ngoặc thứ hai:
		$$ \frac{x}{x+2\sqrt{x}} + \frac{x}{\sqrt{x}+2} = \frac{x}{\sqrt{x}(\sqrt{x}+2)} + \frac{x\sqrt{x}}{\sqrt{x}(\sqrt{x}+2)} = \frac{x+x\sqrt{x}}{\sqrt{x}(\sqrt{x}+2)} = \frac{x(1+\sqrt{x})}{\sqrt{x}(\sqrt{x}+2)} = \frac{\sqrt{x}(1+\sqrt{x})}{\sqrt{x}+2} $$
		Vậy:
		\begin{align*} A &= \frac{\sqrt{x}+1}{(\sqrt{x}+2)^2} : \frac{\sqrt{x}(1+\sqrt{x})}{\sqrt{x}+2} \\ &= \frac{\sqrt{x}+1}{(\sqrt{x}+2)^2} \cdot \frac{\sqrt{x}+2}{\sqrt{x}(1+\sqrt{x})} \end{align*}
		Vì $x>0$, nên $\sqrt{x}+1 \neq 0$ và $\sqrt{x}+2 \neq 0$.
		\begin{align*} A &= \frac{\cancel{(\sqrt{x}+1)}}{(\sqrt{x}+2)^{\cancel{2}}} \cdot \frac{\cancel{(\sqrt{x}+2)}}{\sqrt{x}\cancel{(1+\sqrt{x})}} \\ &= \frac{1}{\sqrt{x}(\sqrt{x}+2)} \end{align*}
		Vậy $A = \frac{1}{x+2\sqrt{x}}$.
		\\ \textbf{Đáp số: $A = \frac{1}{x+2\sqrt{x}}$}
	}
\end{ex}

\begin{ex}
	Biết phương trình $-2x^2+5x-2=0$ có hai nghiệm $x_1; x_2$. Không giải phương trình, hãy tính $M=x_1\left(1-\frac{1}{x_2}\right)+x_2\left(1-\frac{1}{x_1}\right)$.
	\loigiai{
		Phương trình $-2x^2+5x-2=0$ có $a=-2, b=5, c=-2$.
		$\Delta = b^2-4ac = 5^2 - 4(-2)(-2) = 25 - 16 = 9 > 0$.
		Phương trình có hai nghiệm phân biệt $x_1, x_2$.
		Theo hệ thức Vi-ét, ta có:
		$$ \begin{cases} x_1 + x_2 = -\frac{b}{a} = -\frac{5}{-2} = \frac{5}{2} \\ x_1 x_2 = \frac{c}{a} = \frac{-2}{-2} = 1 \end{cases} $$
		Vì $x_1x_2 = 1 \neq 0$, nên $x_1 \neq 0$ và $x_2 \neq 0$.
		Ta biến đổi biểu thức M:
		\begin{align*} M &= x_1\left(\frac{x_2-1}{x_2}\right) + x_2\left(\frac{x_1-1}{x_1}\right) \\ &= \frac{x_1(x_2-1)}{x_2} + \frac{x_2(x_1-1)}{x_1} \\ &= \frac{x_1^2(x_2-1) + x_2^2(x_1-1)}{x_1x_2} \\ &= \frac{x_1^2x_2 - x_1^2 + x_1x_2^2 - x_2^2}{x_1x_2} \\ &= \frac{x_1x_2(x_1+x_2) - (x_1^2+x_2^2)}{x_1x_2} \\ &= \frac{x_1x_2(x_1+x_2) - ((x_1+x_2)^2 - 2x_1x_2)}{x_1x_2} \end{align*}
		Thay giá trị từ hệ thức Vi-ét vào biểu thức M:
		\begin{align*} M &= \frac{1 \cdot \left(\frac{5}{2}\right) - \left(\left(\frac{5}{2}\right)^2 - 2 \cdot 1\right)}{1} \\ &= \frac{5}{2} - \left(\frac{25}{4} - 2\right) \\ &= \frac{5}{2} - \left(\frac{25}{4} - \frac{8}{4}\right) \\ &= \frac{5}{2} - \frac{17}{4} \\ &= \frac{10}{4} - \frac{17}{4} \\ &= -\frac{7}{4} \end{align*}
		Vậy $M = -\frac{7}{4}$.
		\\ \textbf{Đáp số: $M = -\frac{7}{4}$}
	}
\end{ex}

\begin{ex}
	Giải hệ phương trình $\begin{cases} 4x-3y=1 \\ 2x+3y=7 \end{cases}$.
	\loigiai{
		Ta có hệ phương trình:
		$$ \begin{cases} 4x-3y=1 \quad &(1) \\ 2x+3y=7 \quad &(2) \end{cases} $$
		Cộng vế theo vế của phương trình (1) và (2), ta được:
		\begin{align*} (4x-3y) + (2x+3y) &= 1+7 \\ 6x &= 8 \\ x &= \frac{8}{6} \\ x &= \frac{4}{3} \end{align*}
		Thay $x=\frac{4}{3}$ vào phương trình (2):
		\begin{align*} 2\left(\frac{4}{3}\right) + 3y &= 7 \\ \frac{8}{3} + 3y &= 7 \\ 3y &= 7 - \frac{8}{3} \\ 3y &= \frac{21}{3} - \frac{8}{3} \\ 3y &= \frac{13}{3} \\ y &= \frac{13}{9} \end{align*}
		Vậy nghiệm của hệ phương trình là $\left(x;y\right) = \left(\frac{4}{3};\frac{13}{9}\right)$.
		\\ \textbf{Đáp số: $(x;y) = \left(\frac{4}{3};\frac{13}{9}\right)$}
	}
\end{ex}
\begin{ex}
	Bạn Nam rút ngẫu nhiên một tấm thẻ từ trong hộp chứa 30 tấm thẻ ghi các số 11; 12; 13; ...; 40 biết hai thẻ khác nhau thì ghi hai số khác nhau. Tính xác suất của các biến cố A: “Rút được tấm thẻ ghi số chia hết cho 5”.
	\loigiai{
		Các số ghi trên 30 tấm thẻ là các số tự nhiên từ 11 đến 40.
		Không gian mẫu $\Omega = \{11; 12; 13; \dots; 39; 40\}$.
		Số phần tử của không gian mẫu là $n(\Omega) = 40 - 11 + 1 = 30$.
		
		Gọi A là biến cố: “Rút được tấm thẻ ghi số chia hết cho 5”.
 các kết quả thuận lợi cho biến cố A là $A = \{15; 20; 25; 30; 35; 40\}$.
		Số phần tử của biến cố A là $n(A) = 6$.
		
		Xác suất của biến cố A là:
		$P(A) = \frac{n(A)}{n(\Omega)} = \frac{6}{30} = \frac{1}{5}$.
		\\ \textbf{Đáp số: $P(A) = \frac{1}{5}$}
	}
\end{ex}
\begin{ex}
	Ông Hai chia số tiền 500 triệu đồng của mình cho hai khoản đầu tư. Sau một năm, tổng số tiền lãi thu được là 29 triệu đồng. Lãi suất cho khoản đầu tư thứ nhất là 5\% một năm. Lãi suất cho khoản đầu tư thứ hai là 7\% một năm. Tính số tiền ông Hai đầu tư cho mỗi khoản.
	\loigiai{
		Gọi số tiền ông Hai đầu tư cho khoản thứ nhất là $x$ (triệu đồng).
		Gọi số tiền ông Hai đầu tư cho khoản thứ hai là $y$ (triệu đồng).
		Điều kiện: $x > 0, y > 0$.
		
		Tổng số tiền đầu tư là 500 triệu đồng, nên ta có phương trình:
		$x + y = 500 \quad (1)$
		
		Số tiền lãi từ khoản đầu tư thứ nhất sau một năm là: $0,05x$ (triệu đồng).
		Số tiền lãi từ khoản đầu tư thứ hai sau một năm là: $0,07y$ (triệu đồng).
		
		Tổng số tiền lãi thu được là 29 triệu đồng, nên ta có phương trình:
		$0,05x + 0,07y = 29 \quad (2)$
		
		Từ (1) và (2), ta có hệ phương trình:
		$$ \begin{cases} x + y = 500 \\ 0,05x + 0,07y = 29 \end{cases} $$
		Giải hệ phương trình:
		Từ (1) $\Rightarrow y = 500 - x$.
		Thế $y = 500 - x$ vào (2):
		\begin{align*} 0,05x + 0,07(500 - x) &= 29 \\ 0,05x + 35 - 0,07x &= 29 \\ -0,02x &= 29 - 35 \\ -0,02x &= -6 \\ x &= \frac{-6}{-0,02} \\ x &= 300 \end{align*}
		Thay $x = 300$ vào $y = 500 - x$:
		\begin{align*} y &= 500 - 300 \\ y &= 200 \end{align*}
		(Thỏa mãn điều kiện $x > 0, y > 0$).
		
		Kiểm tra lại:
		Tiền lãi từ khoản thứ nhất: $0,05 \times 300 = 15$ triệu đồng.
		Tiền lãi từ khoản thứ hai: $0,07 \times 200 = 14$ triệu đồng.
		Tổng tiền lãi: $15 + 14 = 29$ triệu đồng (đúng).
		
		Vậy, số tiền ông Hai đầu tư cho khoản thứ nhất là 300 triệu đồng và cho khoản thứ hai là 200 triệu đồng.
		\\ \textbf{Đáp số: Khoản thứ nhất: 300 triệu đồng, Khoản thứ hai: 200 triệu đồng}
	}
\end{ex}
\begin{ex}
	Cho đường tròn (O) có hai đường kính AB và MN vuông góc với nhau. Trên tia đối của tia MA lấy điểm C khác điểm M. Kẻ MH vuông góc với BC (H thuộc BC). MB cắt OH tại E.
	\begin{enumerate}
		\item[a)] Chứng minh B, O, M, H cùng thuộc một đường tròn.
		\item[b)] Gọi giao điểm của đường tròn (O) với đường tròn ngoại tiếp tam giác MHC là K. Chứng minh $ME \cdot HB = EB \cdot MH$ và ba điểm C, K, E thẳng hàng.
	\end{enumerate}
\end{ex}


\setcounterref{ex}{0}
